% !TeX root = ../se200_referential_regimes.tex

% =============================================================================
\section{Related Work}
\label{sec:related}
% =============================================================================

This paper is situated at the intersection of formal ontology, provenance, and
classification under disagreement.
It does not propose a new upper ontology.
It derives a lower bound on identity regimes required when a substrate
is optimized for neutral reference under persistent disagreement.

% -----------------------------------------------------------------------------
\subsection{Classification under Disagreement}
\label{subsec:related_classification_disagreement}
% -----------------------------------------------------------------------------

Classification systems embed institutional, social, and evaluative commitments
rather than functioning as neutral descriptions~\citep{bowkerstar1999,haslanger2012}.
In scientific and policy settings, disagreement is often persistent rather than transient,
reflecting incompatible evidentiary, methodological, or evaluative frameworks
rather than temporary empirical uncertainty~\citep{longino1990}.
This motivates representations that remain usable without presupposing interpretive convergence.

This paper adopts that motivation but addresses a narrower structural question:
which identity distinctions become unavoidable when disagreement
is persistent and causal or normative commitment is externalized from the substrate.
The answer is conditional on the neutral-substrate constraint developed
in \emph{Neutral Substrates}~\citep{case2026neutral}
and on the formal setting fixed above.
Under those assumptions, stable reference requires at least the nine core identity regimes derived here.

% -----------------------------------------------------------------------------
\subsection{Identity, Roles, and Ontological Methodology}
\label{subsec:related_identity_roles}
% -----------------------------------------------------------------------------

Formal ontology has long emphasized that categories are individuated
by identity conditions and persistence criteria.
OntoClean clarifies why anti-rigid properties such as roles
cannot supply identity conditions~\citep{guarino1998roles,guarinowelty2002}.
Work on social and institutional entities further develops distinctions among
bearers, roles, social objects, and institutional structures~\citep{masolo2003wonderweb,guizzardi2005}.

This paper builds on those insights but applies them at the substrate level.
Its regimes are not derived from domain taxonomy.
They are derived from transformation behavior over identity bases.
The central move is to make identity operational:
a regime is determined by whether ordinary transformations
preserve, break, or leave identity neutral.

% -----------------------------------------------------------------------------
\subsection{Upper Ontologies}
\label{subsec:related_upper_ontologies}
% -----------------------------------------------------------------------------

Upper ontologies such as BFO, DOLCE, and UFO provide rich accounts of
dependence, endurants and perdurants, social objects, qualities, and
institutional entities~\citep{smithceusters2015bfo,arp2015building,gangemi2002dolce,guizzardi2005}.
They support broad ontology-engineering objectives, including reuse,
domain modeling, semantic integration, and metaphysical clarity.
The present result addresses a narrower design objective:
stable substrate-level reference under persistent disagreement.
It offers a constraint that can be imposed on
implementations when that objective is adopted.

From the perspective of BFO, several of the required carriers have natural
representational starting points.
Occurrences map naturally to occurrents.
Material plain referents map naturally to independent continuants.
Rules, records, roles, dispositions, information artifacts, and institutional
bearers can also be modeled using BFO resources, depending on the intended application.
When a BFO-based implementation is intended to serve as a
neutral substrate under persistent disagreement,
the relevant identity bases can be specified as part of the implementation.
For example, non-spatial normative applicability, record identity, and obligation-bearing capacity
may require dedicated modeling patterns to ensure that the
regime-level distinctions are preserved across transformations.

DOLCE and UFO provide rich resources for social objects,
institutional entities, qualities, roles, and dependence relations.
These resources make them natural candidates for representing
obligation-bearers, rules, and institutional structures.
Here again, the contribution is an additional constraint
imposed when neutrality under persistent
disagreement is adopted as a design objective.
A deployed substrate must make clear when it is
tracking a rule by content rather than structure,
a scope by extension rather than internal organization,
a plain referent by locus rather than object,
and a record as distinct from the thing recorded.
Where those distinctions are handled by
local modeling conventions, role predicates, or application logic,
the implementation may consider whether those conventions remain stable
when shared interpretation cannot be assumed.

These observations are illustrative.
The core is derived from the
transformation analysis of Sections~\ref{sec:algebra}--\ref{sec:bound},
not from any gaps in upper ontologies.
Upper ontologies can realize the regime structure derived here when
their implementations preserve the relevant identity behavior.
The benefit is a reusable constraint for cases where an ontology must support
neutral reference across divergent interpretations.

% -----------------------------------------------------------------------------
\subsection{Provenance, Versioning, and Exchange}
\label{subsec:related_provenance_exchange}
% -----------------------------------------------------------------------------

The transformation basis used in this paper mirrors operations
already tracked by provenance and versioning systems.
PROV represents derivation, attribution, generation, use, and revision
in ways that overlap naturally with re-expression, annotation,
refinement, branching, and provenance extension~\citep{provdm2013}.
This paper adds an identity-level constraint for cases
where provenance and versioning records
must remain referable across incompatible interpretations.

The neutral substrate and its regimes can be stated as a conservative theory
over an exchange formalism such as Common Logic~\citep{isoiec24707}.
Conservativity matters because extensions should add interpretive structure
without changing consequences about substrate-level terms~\citep{hodges1997}.
This supports heterogeneous ontology exchange and pluralistic structuring
in settings where no single interpretive framework can govern all use~\citep{kutz2010}.

Large-scale information infrastructures and digital-twin systems face
related problems of long-lived reference, versioning, and cross-institutional
reuse~\citep{edwards2011,grieves2016}.
The literature emphasizes that infrastructure must remain usable
across evolving communities and interpretations.
This paper provides a constraint that enables stable reference under persistent disagreement.

% -----------------------------------------------------------------------------
\subsection{Causal and Normative Separation}
\label{subsec:related_causal_normative}
% -----------------------------------------------------------------------------

The distinction between descriptive record and causal claim follows a familiar discipline in causal analysis.
Causal claims require additional modeling commitments and should not be confused
with descriptive structure alone~\citep{pearl2009}.
The neutral-substrate constraint applies the same discipline at the substrate boundary.
A substrate may record that a source asserted a causal or normative proposition.
It must not thereby assert the proposition as substrate-level fact.

The regimes derived here make that separation usable under transformation.
Records can persist, be annotated, forked, and extended with provenance
without becoming causal explanations or normative judgments.
Rules can remain referable without the substrate deciding which interpretation governs.
Occurrences can remain referable without the substrate deciding
their causal effects or compliance status.
That separation is essential when identity must remain stable
before interpretation is settled.

% -----------------------------------------------------------------------------
\subsection{FAIR Reuse under Persistent Disagreement}
\label{subsec:related_fair}
% -----------------------------------------------------------------------------

The FAIR principles promote data reuse across communities and over time~\citep{wilkinson2016}.
They do not specify the ontological identity conditions required
when reuse occurs under persistent legal, institutional, or analytic disagreement.
In practice, reuse is often pursued through negotiated
vocabularies, shared schemas, or domain-specific interpretive conventions.
Those strategies are valuable, but they presuppose forms of
convergence that may be unavailable in contested settings.

This result complements FAIR-aligned reuse in contested settings.
Stable reference must be maintained independently of causal or normative agreement.
The nine core regimes name identity distinctions that a neutral substrate
must represent when reuse is expected to continue across divergent interpretations.
